% TT30001C
% Verbrennung, leicht
% 14.0
% 2013-11-30
:ref:`m-verbrennung-leicht`
\begin{itemize}
  -   Selbstschutz!
  -   Kleiderbrand löschen
  -   Vitale Bedrohung einschätzen. |TxBeiVit| |TxMassVitMK|
  -   Kleidung und Schmuck entfernen (sofern nicht  eingeschmolzen!)
  -   Nur wenn \E{unmittelbar nach dem Unfall\footnote{Bereits zwei
  Minuten nach Verbrennungsbeginn und damit bei Eintreffen
  des Rettungsdienstes ist ein
  positiver Effekt der Kühlung nicht mehr zu erwarten. Bei
  mehreren Minuten zurückliegenden Verbrennungen ist eine
  Kühlung nicht mehr sinnvoll.
 :cite:`Asanger2010,AsboeLehrmeinungRd2011-21`\ }}
  damit begonnen werden kann, für max. 10\Min* mit handwarmen
  Wasser die betroffenen Körperregionen spülen;
  
  \EE{Bei Frösteln des Patienten früher abzubrechen: Gefahr der Unterkühlung!}
  -   **Abdeckung**: Sterile, nicht verklebende
  Wundauflagen; evtl. metalliserte Wundauflagen
 :cite:`Helfen2008,RedelsteinerHRNS1`\ ; locker anlegen
  % DISCUSS INHALT: Water-Jel, Stellungsnahme
  %         -   Wenn vorhanden: Water-Jel
  -   Ggfs. Schockbekämpfung
  %   -   Ggfs. Notarzt-Anforderung
\end{itemize}